% L15a strategy deck: the liquid-gamma picking rule.
% One short deck per rule on the strategy menu, read next to the rule's code file
% (strategies/picking/liquid_gamma.jl). Written by make_strategy_decks.py (session scratchpad, 2026-11-27) and then
% kept by hand: edit this file, not the generator. Build with make in this folder: LaTeX runs from strategies/,
% so paths are relative to it.
\documentclass[aspectratio=169,11pt]{beamer}
\usepackage{vnslides}

\newcommand{\repo}{https://github.com/varnerlab/CHEME-5660-CourseRepository-Fall-2026/blob/main/lectures/week-15/L15a}

\title{\texorpdfstring{{\vnheavy L15a}}{L15a}: Liquid Gamma}
\subtitle{Picking Rule, CHEME 5660 Strategy Menu, Cornell University}
\author{Copyright Jeffrey Varner 2026. All Rights Reserved}

\begin{document}

\vntitlepage

\begin{frame}{Disclaimer and Risks}
  \vspace{0.6em}
  \textbf{\ink{Training only.}} This content is for training and informational
  purposes only. The teaching team does not make, give, or endorse any offer or
  solicitation to buy or sell securities or derivative products, or provide
  investment or trading advice or strategy.

  \vspace{1.1em}
  \textbf{\ink{Risk.}} Trading involves risk and is generally not recommended for
  those with limited resources, experience, or a low risk tolerance. Past
  performance does not guarantee future success.

  \vspace{1.1em}
  \textbf{\ink{Responsibility.}} You are responsible for your investment decisions
  based on your financial circumstances, objectives, risk tolerance, and
  liquidity needs.
\end{frame}

\begin{frame}{The Idea}
  \vspace{0.3em}
  The liquid-gamma rule picks the 13 tickers with the highest preference weights among the 100 most traded
  tickers:
  \begin{itemize}\setlength{\itemsep}{0.5em}
    \item \hilite{Most traded}: the 100 tickers with the largest median value traded per day (USD/day) over the last 63 trading days.
    \item \hilite{Reason}: the price impact part of the engine's trading cost grows with the size of an order relative to the ticker's average daily volume. The same order costs less in a heavily traded ticker.
    \item \hilite{Size}: the basket holds $K=13$ tickers at every decision.
  \end{itemize}

  \vspace{0.5em}
  The engine runs the rule at each decision: on the start day and whenever a trigger fires. Between decisions, the basket stays the same.
\end{frame}

\begin{frame}{The Preference Weight}
  \vspace{0.3em}
  At each close $t$, the engine scores every ticker $i$ with its preference weight $\gamma_{i}(t)$:
  {\setlength{\jot}{3pt}
  \begin{align*}
    \tilde{g}_{i,t} &= \alpha_{i}(t)+\beta_{i}(t)\,\tilde{g}_{M,t} && \text{the expected growth rate (1/yr)}\\
    \gamma_{i}(t) &= \tanh\left(h\,\frac{\tilde{g}_{i,t}}{\max\left(|\beta_{i}(t)|,\,\beta_{\text{floor}}\right)^{\xi_{t}}}\right)
      && h=1\text{ year},\ \beta_{\text{floor}}=0.1
  \end{align*}}%
  where the single index model gives $\alpha_{i}(t)$ and $\beta_{i}(t)$, $\tilde{g}_{M,t}$ is the 252-day
  exponential moving average of market growth, and the market state $\xi_{t}$ is positive when the market falls.

  \vspace{0.5em}
  The preference weight lies between $-1$ and $1$, with the sign of the expected growth rate.
\end{frame}

\begin{frame}{How the Market State Tilts the Preference Weight}
  \vspace{0.3em}
  The market state $\xi_{t}$ sets how the market sensitivity scales the expected growth rate:
  \begin{center}
    \includegraphics{slides/figs/Fig-L15a-Deck-Tilt.pdf}
  \end{center}
  At the same positive expected growth rate, the preference weight favors a low market sensitivity when the
  market falls and a high one when the market rises.
\end{frame}

\begin{frame}{The Rule}
  \vspace{0.3em}
  On decision day $t$, let $S_{i}(u)$ (USD/share) and $V_{i}(u)$ (shares) be the closing price and the volume
  of ticker $i$ on day $u$:
  {\setlength{\jot}{3pt}
  \begin{align*}
    D_{i}(t) &= \operatorname{median}_{\,t-62\leq u\leq t}\;S_{i}(u)\,V_{i}(u) && \text{value traded (USD/day)}\\
    \mathcal{L}(t) &= \text{the 100 tickers with the largest }D_{i}(t) && \text{shortlist}\\
    \mathcal{B}(t) &= \text{the 13 tickers in }\mathcal{L}(t)\text{ with the largest }\gamma_{i}(t) && \text{basket}
  \end{align*}}%
  The order of the tickers in the universe breaks ties in both rankings.
\end{frame}

\begin{frame}{The Two Steps of the Rule}
  \vspace{0.3em}
  The rule shortlists the most traded tickers, then ranks the shortlist by preference weight:
  \begin{center}
    \includegraphics{slides/figs/Fig-L15a-Deck-Picks-Liquid-Gamma.pdf}
  \end{center}
  The random-liquid rule draws its 13 tickers from the same shortlist.
\end{frame}

\begin{frame}{Reading the Code}
  \vspace{0.3em}
  The rule's code file is
  \href{\repo/strategies/picking/liquid_gamma.jl}{\texttt{strategies/picking/liquid\_gamma.jl}}:
  \begin{itemize}\setlength{\itemsep}{0.5em}
    \item \texttt{pick\_liquid\_gamma(tickers, gamma, dollar\_volume; count, shortlist)} builds the shortlist, then takes the highest preference weights in it.
    \item \texttt{MyL15aLiquidGammaScreen} is the type the engine calls. It holds the settings \texttt{count~=~13}, \texttt{shortlist~=~100}, and \texttt{window~=~63}.
    \item \texttt{MPC.pick(screen, md, t; $\gamma$\_t, ...)} computes $D_{i}(t)$ with \texttt{MPC.\_liquidity\_scores}, runs the rule with the day's preference weights \texttt{$\gamma$\_t}, and returns the basket's columns in the market data.
  \end{itemize}

  \vspace{0.5em}
  The \href{\repo/strategies/CHEME-5660-L15a-Strategy-Menu-Fall-2026.ipynb}{strategy menu} lists every rule, and the \href{\repo/engine/CHEME-5660-L15a-Trading-Engine-Fall-2026.ipynb}{trading engine} describes the preference weight and the triggers.
\end{frame}

\end{document}
